\chapter{Linux Networking Fundamentals}
\label{chap:linuxnetworking}

From user space, networking starts with the \highlight{Berkeley sockets} API. 
HFT stacks care about \highlight{which} protocol carries market data versus orders, 
and whether a call can \highlight{block} the hot thread (unpredictable wake-ups). 
Multiplexing with \texttt{epoll} is Chapter~\ref{chap:eventdrivenio}; kernel bypass is 
Chapter~\ref{chap:kernelbypass}.
\newline
Do \highlight{not} memorize every flag and \texttt{errno}. Use \texttt{man 2 <call>} and 
\texttt{man 7 socket} / \texttt{man 7 tcp}. Interviewers care that you know which calls 
block, how non-blocking I/O fails with \texttt{EAGAIN}, and latency knobs such as 
\texttt{TCP\char`_NODELAY}.

\section{TCP vs UDP}
\label{sec:tcp-vs-udp}
\begin{itemize}
    \item \highlightbf{TCP} --- connection-oriented, reliable, ordered byte stream. 
          Handshake, retransmission, congestion control. Common for order/execution 
          sessions where loss is unacceptable. Cost: more kernel work and latency variance.
    \item \highlightbf{UDP} --- datagrams, no connection, no delivery guarantee. 
          Common for \highlight{market data} (including multicast feeds): the app handles 
          gaps/seq numbers. Lower overhead; you own reliability.
\end{itemize}
Interview line: pick TCP when you need a reliable session; UDP when you need 
firehose throughput and can tolerate/detect loss yourself.

\section{Which Calls Block?}
\label{sec:socket-blocking-matrix}
By default a new socket is \highlight{blocking}. \texttt{O\char`_NONBLOCK} (via \texttt{fcntl} or 
\texttt{SOCK\char`_NONBLOCK} / \texttt{accept4}) changes that.
\begin{center}
\begin{tabular}{|l|l|p{7.5cm}|}
\hline
\textbf{Call} & \textbf{Default} & \textbf{Notes} \\
\hline
\texttt{socket} & non-blocking & Local kernel alloc only. \\
\hline
\texttt{bind} & non-blocking & Local name binding. \\
\hline
\texttt{listen} & non-blocking & Marks socket passive; sets backlog. \\
\hline
\texttt{accept} & \highlight{blocking} & Sleeps if no pending connection; 
  non-blocking $\rightarrow$ \texttt{EAGAIN}. \\
\hline
\texttt{connect} & \highlight{blocking} & Waits for TCP handshake; 
  non-blocking $\rightarrow$ \texttt{EINPROGRESS}. \\
\hline
\texttt{send} & \highlight{blocking} & Sleeps if send buffer full; may send \highlight{partially}. \\
\hline
\texttt{recv} & \highlight{blocking} & Sleeps if no data; return 0 = peer closed cleanly. \\
\hline
\texttt{close} & usually not & Returns quickly; kernel finishes teardown async. 
  Can block with \texttt{SO\char`_LINGER}. \\
\hline
\end{tabular}
\end{center}

\section{Socket Lifecycle}
\label{sec:socket-lifecycle}
A socket is an endpoint (an \texttt{fd}). Typical TCP server path: 
\texttt{socket} $\rightarrow$ \texttt{setsockopt} $\rightarrow$ \texttt{bind} $\rightarrow$ 
\texttt{listen} $\rightarrow$ \texttt{accept} loop $\rightarrow$ \texttt{recv}/\texttt{send} 
$\rightarrow$ \texttt{close}. 
Client path: \texttt{socket} $\rightarrow$ \texttt{connect} $\rightarrow$ I/O $\rightarrow$ \texttt{close}. 
UDP often skips \texttt{listen}/\texttt{accept} and uses \texttt{sendto}/\texttt{recvfrom} 
(or \texttt{connect} on UDP to fix a peer).

\begin{figure}[H]
    % TCP socket lifecycle: server (left) vs client (right)
\centering
\begin{tikzpicture}[
    font=\sffamily,
    box/.style={
        draw=black,
        thick,
        rectangle,
        minimum width=2.6cm,
        minimum height=0.72cm,
        align=center,
        text=oblivionkeyword
    },
    title/.style={
        text=oblivionkeyword,
        font=\sffamily\bfseries\large
    },
    arr/.style={-{Stealth[length=2.2mm]}, thick, black}
]
    % --- titles ---
    \node[title] (srvtitle) at (0,0) {Server};
    \node[title] (cltitle) at (5.2,0) {Client};

    % --- server column ---
    \node[box] (ssocket) at (0,-1.0) {Socket};
    \node[box] (ssetsockopt) at (0,-2.1) {setsockopt};
    \node[box] (sbind) at (0,-3.2) {Bind};
    \node[box] (slisten) at (0,-4.3) {Listen};
    \node[box] (saccept) at (0,-5.4) {Accept};
    \node[box] (ssendrecv) at (0,-6.5) {Send/Recv};

    \draw[arr] (ssocket) -- (ssetsockopt);
    \draw[arr] (ssetsockopt) -- (sbind);
    \draw[arr] (sbind) -- (slisten);
    \draw[arr] (slisten) -- (saccept);
    \draw[arr] (saccept) -- (ssendrecv);

    % --- client column ---
    \node[box] (csocket) at (5.2,-1.0) {Socket};
    % Connect between Listen and Accept
    \node[box] (cconnect) at (5.2,-4.85) {Connect};
    \node[box] (csendrecv) at (5.2,-6.5) {send/Recv};

    \draw[arr] (csocket) -- (cconnect);
    \draw[arr] (cconnect) -- (csendrecv);

    % --- interactions ---
    \draw[arr, bend right=28] (cconnect.north west) to (slisten.east);
    \draw[{Stealth[length=2.2mm]}-{Stealth[length=2.2mm]}, thick, black]
        (ssendrecv) -- (csendrecv);
\end{tikzpicture}

    \caption{TCP socket lifecycle: server setup versus client \texttt{connect}, then bidirectional I/O.}
    \label{fig:socket-lifecycle}
\end{figure}

\subsection{\texttt{socket()}}
\label{subsec:socket}
Creates an endpoint. Example: \texttt{socket(AF\char`_INET, SOCK\char`_STREAM, 0)} for TCP IPv4; 
\texttt{SOCK\char`_DGRAM} for UDP. Returns a file descriptor or \texttt{-1} on error 
(e.g.\ \texttt{EMFILE} if the process is out of fds). Not a blocking wait on the network.

\subsection{\texttt{bind()}}
\label{subsec:bind}
Assigns a local address (IP + port). Servers bind before listening; 
\texttt{INADDR\char`_ANY} means ``all interfaces.'' Not a network wait.
\newline
\highlightbf{\texttt{EADDRINUSE} / \texttt{TIME\char`_WAIT}:} after a restart, \texttt{bind} often fails 
because the port is still reserved while TCP finishes delayed segments. Set 
\texttt{SO\char`_REUSEADDR} with \texttt{setsockopt} \highlight{before} \texttt{bind} so the address 
can be reused. This is a standard interview pitfall.

\subsection{\texttt{listen()}}
\label{subsec:listen}
Marks a TCP socket as passive. The \texttt{backlog} limits the queue of connections that 
finished the handshake but are not yet taken by \texttt{accept}. If the queue is full, 
clients may see \texttt{ECONNREFUSED}. System cap: \texttt{/proc/sys/net/core/somaxconn}.

\subsection{\texttt{accept()}}
\label{subsec:accept}
Dequeues a completed connection and returns a \highlight{new} fd for that client. 
Blocking \texttt{accept} sleeps until a client arrives. Non-blocking: \texttt{-1} with 
\texttt{EAGAIN}/\texttt{EWOULDBLOCK}. Optional Linux \texttt{accept4} can set 
\texttt{SOCK\char`_NONBLOCK} / \texttt{SOCK\char`_CLOEXEC} atomically on the new fd.

\subsection{\texttt{connect()}}
\label{subsec:connect}
TCP client: starts the three-way handshake. Blocking \texttt{connect} may wait a long time 
on failure/timeout. Non-blocking: returns \texttt{-1} / \texttt{EINPROGRESS}; completion is 
detected by polling for \highlight{writability} (\texttt{epoll}/\texttt{poll}), then checking 
socket error via \texttt{getsockopt(SO\char`_ERROR)}. If the client did not \texttt{bind}, the 
kernel picks an ephemeral local port.

\subsection{\texttt{send()}}
\label{subsec:send}
Copies bytes into the kernel send buffer (TCP stream or connected UDP). May accept 
\highlight{fewer} bytes than requested --- loop until all data is queued. Blocking sleeps when 
the buffer is full; non-blocking yields \texttt{EAGAIN}.
\newline
\highlightbf{\texttt{MSG\char`_NOSIGNAL}:} writing to a closed peer can raise \texttt{SIGPIPE} and 
\highlight{kill} the process. Pass \texttt{MSG\char`_NOSIGNAL} so failure becomes \texttt{-1} / 
\texttt{EPIPE} instead --- critical for trading processes that must not die on a dropped session.

\subsection{\texttt{recv()}}
\label{subsec:recv}
\begin{itemize}
    \item \texttt{> 0} --- bytes read (may be partial).
    \item \texttt{0} --- peer closed cleanly (\texttt{close}/\texttt{shutdown}).
    \item \texttt{-1} --- error (\texttt{EAGAIN} if non-blocking and empty; \texttt{ECONNRESET} if peer aborted).
\end{itemize}

\subsection{\texttt{close()}}
\label{subsec:close}
Drops one reference on the fd. TCP teardown usually continues in the kernel without 
blocking your thread. With \texttt{SO\char`_LINGER}, \texttt{close} can block until data is flushed 
or a timeout expires. After \texttt{fork}, the connection stays up until \highlight{all} processes 
that hold the fd call \texttt{close}; use \texttt{shutdown} when you need to signal end-of-stream 
regardless of refcount.

\begin{lstlisting}[language=C++]
#include <sys/socket.h>
#include <netinet/in.h>
#include <netinet/tcp.h>
#include <unistd.h>

// Minimal TCP listen socket (error checks omitted for brevity)
int make_tcp_listener(uint16_t port)
{
    int fd = socket(AF_INET, SOCK_STREAM, 0);
    int yes = 1;
    setsockopt(fd, SOL_SOCKET, SO_REUSEADDR, &yes, sizeof(yes));

    sockaddr_in addr{};
    addr.sin_family = AF_INET;
    addr.sin_addr.s_addr = htonl(INADDR_ANY);
    addr.sin_port = htons(port);

    bind(fd, reinterpret_cast<sockaddr*>(&addr), sizeof(addr));
    listen(fd, 128);
    return fd;
}
\end{lstlisting}

\section{Blocking vs Non-Blocking I/O}
\label{sec:blocking-vs-nonblocking-io}
This is the HFT-critical distinction in this chapter.

\subsection{Blocking Sockets}
\label{subsec:blocking-sockets}
By default, \texttt{recv}/\texttt{accept}/\texttt{connect}/\texttt{send} may \highlight{sleep} in the kernel. 
The thread is descheduled; wake-up latency is unpredictable 
(Chapter~\ref{chap:linuxprocesses}). Fine for tools; dangerous on a pinned strategy core.

\subsection{Non-Blocking Sockets}
\label{subsec:nonblocking-sockets}
Operations return immediately. If they cannot complete, you get \texttt{EAGAIN} / 
\texttt{EWOULDBLOCK} (or \texttt{EINPROGRESS} for \texttt{connect}) and try later --- typically after 
\texttt{epoll} says the fd is ready (Chapter~\ref{chap:eventdrivenio}), or in a busy-poll loop 
(Chapter~\ref{chap:lowlatency}).

\subsection{\texttt{O\_NONBLOCK}}
\label{subsec:o-nonblock}
Set the file status flag so all I/O on that fd is non-blocking:
\begin{lstlisting}[language=C++]
#include <fcntl.h>

bool set_nonblocking(int fd)
{
    int flags = fcntl(fd, F_GETFL, 0);
    if (flags < 0) return false;
    return fcntl(fd, F_SETFL, flags | O_NONBLOCK) == 0;
}
\end{lstlisting}

\subsection{\texttt{MSG\_DONTWAIT}}
\label{subsec:msg-dontwait}
Per-call non-blocking for \texttt{recv}/\texttt{send} without changing the socket's default mode: 
pass \texttt{MSG\char`_DONTWAIT} in the flags argument. Useful when only some reads must not block.

\section{Latency Knobs}
\label{sec:tcp-latency-knobs}

\subsection{\texttt{TCP\_NODELAY}}
\label{subsec:tcp-nodelay}
By default TCP may use \highlight{Nagle's algorithm}: small writes can be delayed to coalesce 
packets (better bandwidth, worse latency). In HFT that delay is unacceptable.
\begin{lstlisting}[language=C++]
int one = 1;
setsockopt(fd, IPPROTO_TCP, TCP_NODELAY, &one, sizeof(one));
\end{lstlisting}
Disabling Nagle sends ASAP at the cost of more packets. Classic interview question.

\subsection{\texttt{MSG\_NOSIGNAL}}
\label{subsec:msg-nosignal}
See \texttt{send} above: avoid \texttt{SIGPIPE} killing the process; prefer handling \texttt{EPIPE}.

\section{Common Errors}
\label{sec:network-common-errors}
Always check return values; use \texttt{perror} / \texttt{strerror(errno)} while debugging. 
Know these by \highlight{name and meaning}, not numeric codes:

\subsection{\texttt{EAGAIN} / \texttt{EWOULDBLOCK}}
\label{subsec:eagain}
\label{subsec:ewouldblock}
``Try again'' on a non-blocking socket (no data / no send buffer space). On Linux they are 
the same value for sockets. Normal in event-driven loops --- not a fatal error.

\subsection{\texttt{EINPROGRESS}}
\label{subsec:einprogress}
Non-blocking \texttt{connect} started; wait for completion via I/O multiplexing.

\subsection{\texttt{EADDRINUSE}}
\label{subsec:eaddrinuse}
Address/port still busy (often \texttt{TIME\char`_WAIT}). Use \texttt{SO\char`_REUSEADDR} before \texttt{bind}.

\subsection{\texttt{ECONNREFUSED}}
\label{subsec:econnrefused}
Peer actively refused (nothing listening, or listen backlog full).

\subsection{\texttt{ETIMEDOUT}}
\label{subsec:etimedout}
No timely response (path/firewall/server down) --- different from an active refuse.

\subsection{\texttt{EPIPE} / \texttt{ECONNRESET}}
\label{subsec:epipe}
Peer gone. With \texttt{MSG\char`_NOSIGNAL}, broken writes surface as \texttt{EPIPE} instead of 
\texttt{SIGPIPE}. \texttt{ECONNRESET} often means an abrupt abort.
